\subsection{The response-only objective}

Let $y_{b,0},\ldots,y_{b,T}$ be a formatted sequence and let
$m_{b,t}\in\{0,1\}$ indicate whether target token $y_{b,t+1}$ belongs to a supervised
assistant span. With $M=\sum_{b,t}m_{b,t}$, response-only SFT minimizes
\begin{equation}
  \mathcal{L}_{\mathrm{SFT}}(\theta)
  =-\frac{1}{M}\sum_{b=1}^{B}\sum_{t=0}^{T-1}
    m_{b,t}\log p_\theta(y_{b,t+1}\mid y_{b,\leq t}).
  \label{eq:sft-loss}
\end{equation}
This is ordinary next-token cross-entropy with a target-selection mask. If the model
assigns probability $0.4$ to one demonstrated token, that position contributes
$-\log(0.4)\approx0.916$ nats. Lowering the loss moves probability toward observed
tokens and away from their competitors at the same prefix.

The denominator $M$ normalizes the gradient scale; it does not save a Transformer
forward. All $BT$ input positions still pass through the blocks. The mask decides which
positions directly contribute to the scalar loss after their logits have been computed.

Equation~\ref{eq:sft-loss} gives every supervised token equal weight, so a response with
1000 tokens contributes roughly ten times as much total loss as one with 100 tokens. An
alternative first averages within each example and then averages examples:
\begin{equation}
  \mathcal{L}_{\mathrm{example}}
  =-\frac{1}{B}\sum_{b=1}^{B}\frac{1}{M_b}
    \sum_{t=0}^{T-1}m_{b,t}
    \log p_\theta(y_{b,t+1}\mid y_{b,\leq t}),
\end{equation}
where $M_b=\sum_t m_{b,t}$. Neither reduction is universally correct; they express
different choices about whether tokens or demonstrations receive equal total weight.

\paragraph{What $M$ means in current libraries.}
$M$ is not the configured maximum sequence length. In the current Hugging Face causal-LM
loss, labels are shifted once, positions with the ignore index $-100$ are excluded, and
the summed loss is divided by the number of remaining prediction targets across the
effective accumulated batch \citep{huggingface2026gradientaccumulation,
huggingface2026transformersloss}. TRL constructs the ignored positions from completion or
assistant masks; its \texttt{max\_length} instead controls truncation and, when packing is
enabled, the packed block length \citep{huggingface2026trlsft}. Torchtune follows the same
basic convention by counting labels unequal to its ignore index and dividing accumulated
loss by that valid-token count \citep{pytorch2026torchtuneloss}.

Consequently, a padded batch with tensor shape $[B,T]$ still has
\begin{equation}
  M=\sum_{b=1}^{B}\sum_{t=0}^{T-1}
    \mathbf{1}\{\text{shifted label}_{b,t}\neq -100\}\leq BT.
\end{equation}
If full-sequence loss and packing make nearly every position valid, $M$ can be close to
$B$ times the packed length. Under response-only SFT it may be much smaller and varies
with the number of assistant tokens, regardless of how much padding or masked prompt
context occupies the tensor.
